Indeed, $H=N^0$ will satisfy the desired conditions; on the other
hand, if $H$ satisfies them, then (since $H$ normalizes its Lie
algebra $\mathfrak h$) one has $H\subset N$, so, since this is an
inclusion of smooth groups with the same Lie algebra, with $H$
connected, one will have $H=N^0$. For the identity
$\Norm_G(H)=\Norm_G(\mathfrak h)$, one may suppose $k$ algebraically
closed; then from what one has just seen it follows immediately
that the points of the two groups with values in $k$ are the same;
on the other hand the inclusions $H\subset\Norm_G(H)\subset N$
show that $\Norm_G(H)$ and $N$ have the same Lie algebra, hence
they are identical.

\begin{corollary}{5.4}\label{XIII.5.4}
Under the equivalent conditions (i) to (iv), let $a\in\mathfrak h$.
Then the following conditions are equivalent, and hold if $a$ is
regular in $\mathfrak g$:
\begin{enumeratei}
\item $\varphi$ is smooth at $(e,a)$.
\item $M_a=\Transp_G(a,\mathfrak h)$ is smooth at $e$, and
$\dim_e(M_a)=\operatorname{rank}_k\mathfrak h$.
\item $\ad(a)_{\mathfrak g/\mathfrak h}$ is injective (or again,
bijective).
\end{enumeratei}
When one is under the equivalent conditions (xi) to (xiii), let
$H$ be the algebraic subgroup of $G$ considered in 5.3. Then the
preceding conditions are also equivalent to the following conditions:
\begin{enumerate}[label=\textup{(\roman*)},start=4]
\item $\psi$ is \'etale at $(\bar e,a)$.
\item Denoting by $M_a^0$ the connected component of $e$ in the
transporter $M_a$ of $a$ into $\mathfrak h$, equipped with the
structure induced by $M_a$, one has
\[
 H=M_a^0.
\]
\end{enumerate}
\end{corollary}
Clearly (i) $\Rightarrow$ (ii), since $M_a$ is isomorphic to the
fiber $\varphi^{-1}(a)$, the point $e$ corresponding to $(e,a)$,
and one has (ii) $\Rightarrow$ (i), for (ii) implies that $\varphi$
is ``equidimensional'' at $(e,a)$ (i.e. the dimension of the fiber
through this point is that of the generic fiber), which implies
($G\times\mathrm W(\mathfrak h)$ and $\mathrm W(\mathfrak g)$ being
regular) that it is flat at $(e,a)$, hence smooth since its fiber
is so at this point. The equivalence of (i) and (iii) was seen in
the proof of 5.1 as resulting from the simple calculation of the
tangent map. Moreover, one has seen in 4.8 b) that ``$a$ regular
in $\mathfrak g$'' $\Rightarrow$ (iii). Under conditions (xi) to
(xiii), since $q:G\times\mathrm W(\mathfrak h)\to X$ is smooth
($N$ being smooth), it follows that (i) is equivalent to $\psi$
smooth at $(\bar e,a)$, and since $\psi$ is generically \'etale,
this is equivalent to (iv). Finally, as was noted at the end of
the proof of 1.1, (iv) implies that $N$ is the prescheme induced
on $M_a$ by an open and closed part of $M_a$, whence (v); finally
(v) $\Rightarrow$ (ii) trivially (or again (v) $\Rightarrow$ (iv)
by 1.1, for $\psi$ being dominant and $\mathrm W(\mathfrak g)$
normal, ``unramified'' is here equivalent to ``\'etale''); this
completes the proof of 5.4.

\begin{corollary}{5.5}\label{XIII.5.5}
Let $G$ be a smooth algebraic group over a field $k$, $H$ a smooth
algebraic subgroup such that its Lie algebra $\mathfrak h$ contains
(at least over a suitable extension of $k$) a Cartan subalgebra of
the Lie algebra $\mathfrak g$ of $G$. Let $K$ be a connected
algebraic subgroup of $G$ (not necessarily smooth), with Lie algebra
$\mathfrak k$; suppose that $\mathfrak k$ contains a regular element
$a$ of $\mathfrak g$ (at least over a suitable extension of $k$).
Then $H$ contains $K$ if and only if $\mathfrak h$ contains
$\mathfrak k$.
\end{corollary}
Indeed, by 5.4 (iii) $\Rightarrow$ (v) one has $H=M_a^0$ (N.B.
of course, this relation being invariant under extension of the
base field, it is valid without the hypothesis that the latter
be infinite!); on the other hand $\mathfrak k\subset\mathfrak h$
clearly implies $K\subset M_a$, so, since $K$ is connected,
$K\subset M_a^0$, whence $K\subset H$. \hfill Q.E.D.
% typo?: kept as printed: H=M_a^0 although H was not assumed connected.

\begin{corollary}{5.6}\label{XIII.5.6}
Let $G,H$ be as in 5.5 and let $K$ be an algebraic subgroup of
$G$; one supposes $H$ connected and $K$ smooth. Then $K$ contains
$H$ if and only if $\mathfrak k$ contains $\mathfrak h$.
\end{corollary}
Indeed, if $\mathfrak k\supset\mathfrak h$, then $K$ satisfies
the hypothesis considered in 5.5 for $H$; on the other hand $H$
clearly satisfies the condition considered for $K$ in 5.5. The
conclusion then follows from 5.5.

\begin{corollary}{5.7}\label{XIII.5.7}
Let $\mathfrak g$ be the Lie algebra of a smooth algebraic group
$G$ over a field $k$. Then:

\noindent a) Let $\mathfrak d$ be a Lie subalgebra of $\mathfrak g$.
For $\mathfrak d$ to be a Cartan subalgebra, it is necessary and
sufficient that $\mathfrak d$ be nilpotent and equal to its own
normalizer.

\noindent b) Let $a$ be an element of $\mathfrak g$; for $a$ to
be regular, it is necessary and sufficient that
$\operatorname{Nil}(a,\mathfrak g)$ be nilpotent.
\end{corollary}
Extending the base field, one may suppose $k$ infinite. Taking
4.4 into account, one is reduced for a) to proving that if
$\mathfrak d$ is nilpotent and contains an element $a$ such that
$\ad(a)_{\mathfrak g/\mathfrak d}$ is injective, then
$\mathfrak d$ is a Cartan subalgebra. Now by 5.1 (ii)
$\Rightarrow$ (i), $\mathfrak d$ contains a Cartan subalgebra
$\mathfrak d'=\operatorname{Nil}(a,\mathfrak g)$, and by 4.3 one
concludes from the fact that $\mathfrak d$ is nilpotent that
$\mathfrak d=\mathfrak d'$. To prove b), one notes that
$\operatorname{Nil}(a,\mathfrak g)$ is a Cartan subalgebra of
$\mathfrak g$ by a), hence $a$ is regular.

\begin{corollary}{5.8}\label{XIII.5.8}
Let $G$ be a smooth algebraic group over a field $k$, $H$ a smooth
algebraic subgroup, and $\mathfrak g,\mathfrak h$ the Lie algebras;
suppose that after a suitable extension of the base field,
$\mathfrak h$ contains a Cartan subalgebra of $\mathfrak g$.
Let $\mathfrak d$ be a subalgebra of $\mathfrak h$; then it is
a Cartan subalgebra of $\mathfrak h$ if and only if it is a
Cartan subalgebra of $\mathfrak g$.
\end{corollary}
Taking 4.8 c) into account, it remains to show that if
$\mathfrak d$ is a Cartan subalgebra of $\mathfrak h$, it is a
Cartan subalgebra of $\mathfrak g$, and for this one is reduced
to showing that $\mathfrak d$ contains an element $a$ regular
in $\mathfrak g$, supposing (as one may) $k$ algebraically
closed. But since there is a dense open subset of
$\mathfrak h$ consisting of regular points of $\mathfrak g$,
our assertion follows from 5.1 (i) $\Rightarrow$ (vii) applied
to $(\mathfrak h,\mathfrak d)$.

\sgasection{6}{Cartan subalgebras and subgroups of type (C), relative to a smooth algebraic group}
For simplicity, we confine ourselves in the following theorem
to the case of an algebraically closed base field (the case of
an arbitrary base prescheme being treated in the next expos\'e):

\begin{theorem}{6.1}\label{XIII.6.1}
Let $G$ be a smooth algebraic group over an algebraically closed
field $k$, and $\mathfrak g$ its Lie algebra. Then:

\noindent a) The Cartan subalgebras of $\mathfrak g$ are conjugate.

\noindent b) Let $\mathfrak d$ be a Cartan subalgebra of
$\mathfrak g$. Then its normalizer $N$ in $G$ is smooth, and
$D=N^0$ is the only smooth connected subgroup of $G$ whose
Lie algebra is $\mathfrak d$. One has
\[
 \Norm_G(\mathfrak d)=\Norm_G(D)=N,
 \quad\text{hence }D=\Norm_G(D)^0.
\]
\noindent c) With $\mathfrak d$ as in b), set, as in \No 5,
$X=G\times^N\mathrm W(\mathfrak d)$, and consider the
canonical morphism
\[
 \psi:X\to\mathrm W(\mathfrak g),
\]
(whose fiber at $a\in\mathfrak g$ has as its points with
values in $k$ the set of Cartan subalgebras of $\mathfrak g$
which contain $a$). Then $\psi$ is a birational morphism.

\noindent d) With the notation of c), let $U$ be the largest
open subset of $\mathrm W(\mathfrak g)$ such that $\psi$
induces an isomorphism
\[
 \psi^{-1}(U)\isomto U.
\]
Then for $a\in\mathfrak g$, the following conditions are equivalent:
\begin{enumeratei}
\item $a\in U(k)$.
\item $a$ is contained in one and only one Cartan subalgebra
of $\mathfrak g$.
\item The set of Cartan subalgebras of $\mathfrak g$ containing
$a$ is finite and nonempty.
\item The fiber $\psi^{-1}(a)$ has an isolated point.
\item (If $a\in\mathfrak d$) The morphism $\psi$ is \'etale
(or merely: quasi-finite) at the point $(\bar e,a)$.
\item $a$ is a regular element of $\mathfrak g$.
\item With the notation of 5.4 (iii), one has $N=M_a$.
\end{enumeratei}
\end{theorem}
\begin{proof}
Let us apply 5.1 when $\mathfrak h$ is a Cartan subalgebra
$\mathfrak d$ of $\mathfrak g$; let us show that the strongest
conditions (xi) to (xiii) then hold. This is clear, either
in form (xiii), taking 4.7 into account (which implies that
for a regular element $a$ of $\mathfrak g$ contained in
$\mathfrak d$, one has $N=M_a$), or in form (xi) $=$ (x) $+$
(i), for condition (i) is trivial and condition (x) follows
from the fact that a regular point of $\mathfrak g$ is
contained in only one Cartan subalgebra of $\mathfrak g$
(4.6 a)), and a fortiori in only one conjugate of
$\mathfrak d$. Then b) follows from 5.3, and c) follows
from the fact that $\psi$ is generically \'etale and that
a sufficiently general point (more precisely, a regular
point) of $\mathfrak g$ is contained in only one Cartan
subalgebra of $\mathfrak g$. Under these conditions, the
equivalence of conditions (i) to (v) on $a$ is an immediate
consequence of Zariski's Main Theorem for the separated
birational morphism $\psi$, taking into account that
$\mathrm W(\mathfrak g)$ is normal and $X$ integral.
The equivalence of (v) and (vi) is a particular case of
5.4 (iii) $\Leftrightarrow$ (iv) (reducing to the case where
$a\in\mathfrak d$ by transforming $a$ by a suitable element
$g\in G(k)$), taking 4.6 b) into account. Moreover, by 5.4,
(v) and (vi) are also equivalent to $N\supset M_a^0$, and
by 4.7, already invoked, this even implies $N=M$.
% typo?: kept as printed: M rather than M_a.
This proves d).

Of course, b), c), and the equivalence of (i), (iv), (v),
(vi), (vii) remain valid over an arbitrary base field.
Let us now prove a), using the fact that $k$ is algebraically
closed. By 5.1 (i) $\Rightarrow$ (vii),
$\psi:X\to\mathrm W(\mathfrak g)$ is dominant, hence there
exists a dense open subset $V$ of $\mathrm W(\mathfrak g)$
such that every $a\in V(k)$ is the image of an element
of $X(k)$, i.e. is contained in a conjugate of
$\mathfrak d$. Applying this result to a second Cartan
subalgebra $\mathfrak d'$ of $\mathfrak g$, one sees
that one may take $V$ such that every element of $V(k)$
is conjugate to an element of $\mathfrak d$ and to an
element of $\mathfrak d'$. Taking a regular element
in $V(k)$, it follows that there is a conjugate
$\mathfrak d''$ of $\mathfrak d'$ which contains an
element of $\mathfrak d$ regular in $\mathfrak g$,
and hence is equal to $\mathfrak d$ by 4.6 a).
This completes the proof of 6.1.
\end{proof}

\begin{definition}{6.2}\label{XIII.6.2}
Let $G$ be a smooth algebraic group over a field $k$.
One calls a \emph{subgroup of type (C)} of $G$ any
smooth connected algebraic subgroup whose Lie algebra
is a Cartan subalgebra of $\mathfrak g$. One calls the
\emph{infinitesimal rank} of $G$ the nilpotent rank of
its Lie algebra $\mathfrak g$, equal to the dimension
of every subgroup of type (C) of $G$.
\end{definition}
From 6.1 b) one concludes at once:

\begin{corollary}{6.3}\label{XIII.6.3}
Under the conditions of 6.2, the map
$D\mto\mathfrak d=\Lie D$ establishes a one-to-one
correspondence between subgroups of type (C) of $G$
and Cartan subalgebras $\mathfrak d$ of $\mathfrak g$.
If $D$ is a subgroup of type (C) of $G$, $D$ is its
own connected normalizer:
\[
 D=\Norm_G(D)^0.
\]
\end{corollary}
Combining 6.1 (a) and 6.3, one finds:

\begin{corollary}{6.4}\label{XIII.6.4}
If $k$ is algebraically closed, the subgroups of type
(C) of $G$ are conjugate to one another.
\end{corollary}

\begin{corollary}{6.5}\label{XIII.6.5}
Let $G$ be a smooth algebraic group over algebraically
closed $k$, $H$ a smooth algebraic subgroup of $G$,
and $\mathfrak g$ and $\mathfrak h$ the Lie algebras.
For $\mathfrak h$ to contain a Cartan subalgebra of
$\mathfrak g$, it is necessary and sufficient that $H$
contain a subgroup $D$ of type (C) of $G$.
\end{corollary}
This is clearly sufficient, and it is also necessary,
for $H\supset D$ holds if and only if
$\mathfrak h\supset\mathfrak d$, by 5.5.

\begin{remarks}{6.6}\label{XIII.6.6}
\noindent a) One will note that the connected subgroups
$H$ of $G$ described in 6.5 correspond one-to-one to
the Lie subalgebras $\mathfrak h$ of $\mathfrak g$
which satisfy the strongest condition (xii) of 5.1
(and by 5.5, the inclusion relations between such
subgroups can already be recognized on their Lie algebras).

\noindent b) Still suppose $k$ algebraically closed,
and let $D$ be a subgroup of type (C) of $G$; then
it is easy to show that $D$ contains a Cartan subgroup
$C$ of $G$: indeed, let $V=\mathfrak g/\mathfrak d$;
then for a ``general'' element $a$ of $\mathfrak d$,
$\ad(a)$ acting on $V$ is injective (it suffices that
$a$ be regular in $\mathfrak g$), whence one easily
concludes that for ``general'' $g\in D(k)$,
$\Ad(g)$ acting on $V$ has no nonzero invariant,
which allows one to apply 2.1 (vii) $\Rightarrow$ (i).
In fact, we shall see in the next expos\'e a more
precise result: every Cartan subgroup $C$ of $G$
is contained in one and only one subgroup of type
(C) of $G$.

\noindent c) One should take care not to confuse
Cartan subgroups and subgroups of type (C) in general:
the subgroups of type (C) of $G$ are Cartan subgroups
if and only if they are nilpotent (Cartan subgroups
being indeed maximal nilpotent subgroups); now it
may happen that $\mathfrak g$ is nilpotent without
$G$ being so (example: $\SL(2)_k$ for $k$ of
characteristic $2$); then the Cartan subalgebras
of $\mathfrak g$ are identical to $\mathfrak g$,
i.e. $G$ is a subgroup of type (C) of itself if
$G$ is connected, but it is not a Cartan subgroup
of $G$! On the other hand, we shall see in XIV
that if $G$ is an adjoint semisimple group, then
its subgroups of type (C) are its Cartan subgroups.
Likewise, in characteristic $0$, without restriction
on the smooth algebraic group $G$ over $k$, the
same conclusion holds: this follows at once from
the fact that in characteristic $0$, a connected
algebraic group is nilpotent if (and only if) its
Lie algebra is. This follows at once from the fact
that in characteristic zero, the center $Z$ of
the connected algebraic group $G$ has as Lie algebra
the center of $\mathfrak g$ (for one obtains a priori
the space of invariants of $\mathfrak g$ under the
adjoint representation of $G$; now $G$ being connected
and $k$ of characteristic zero, these are also the
``invariants'' in the infinitesimal sense of the
adjoint representation of $\mathfrak g$, i.e. the
center of $\mathfrak g$), and that by Cartier's
theorem (cf. VI\textsubscript{B} 1.6.1) $Z$ is smooth.
\end{remarks}
